\documentclass[12pt,letterpaper]{article}

% ---------------------------------------------------------------------------
% Packages
% ---------------------------------------------------------------------------
\usepackage[margin=0.9in]{geometry}
\usepackage{newtxtext,newtxmath}
\usepackage{setspace}
\usepackage{titlesec}
\usepackage{booktabs}
\usepackage{array}
\usepackage{tabularx}
\usepackage{enumitem}
\usepackage{parskip}
\usepackage{microtype}
\usepackage[hidelinks]{hyperref}
\usepackage{fancyhdr}
\usepackage{xcolor}

% ---------------------------------------------------------------------------
% Page style
% ---------------------------------------------------------------------------
\pagestyle{fancy}
\fancyhf{}
\rhead{\small SecureGate Explanation Guide}
\lhead{\small CIS 5690}
\rfoot{\small Page \thepage}
\renewcommand{\headrulewidth}{0.4pt}

% ---------------------------------------------------------------------------
% Section formatting
% ---------------------------------------------------------------------------
\titleformat{\section}{\large\bfseries}{}{0em}{}[\titlerule]
\titleformat{\subsection}{\normalsize\bfseries}{}{0em}{}
\titleformat{\subsubsection}{\normalsize\bfseries}{}{0em}{}
\titlespacing*{\section}{0pt}{12pt}{6pt}
\titlespacing*{\subsection}{0pt}{8pt}{4pt}

% ---------------------------------------------------------------------------
% Helpers
% ---------------------------------------------------------------------------
\newcolumntype{L}[1]{>{\raggedright\arraybackslash}p{#1}}
\newcolumntype{Y}{>{\raggedright\arraybackslash}X}

\setlength{\parindent}{0pt}
\setlength{\parskip}{4pt}
\setlength{\headheight}{14pt}
\setlength{\emergencystretch}{1.5em}
\setlist[itemize]{noitemsep, topsep=2pt, leftmargin=*, parsep=2pt}
\setlist[enumerate]{noitemsep, topsep=2pt, leftmargin=*, parsep=2pt}
\setstretch{1.10}

\newcommand{\simple}[1]{\textbf{Simple explanation:} #1}
\newcommand{\technical}[1]{\textbf{Technical explanation:} #1}
\newcommand{\say}[1]{\textbf{What you can say:} ``#1''}

\begin{document}

% ---------------------------------------------------------------------------
% Title block
% ---------------------------------------------------------------------------
\begin{center}
    {\LARGE \textbf{SecureGate}} \\[4pt]
    {\large Complete Explanation Guide for Presentation and Viva} \\[6pt]
    {\normalsize Based on the structured report in \texttt{main.tex}} \\[2pt]
\end{center}
\vspace{2pt}
\hrule
\vspace{10pt}

% ---------------------------------------------------------------------------
\section{Purpose of This Document}

This document is written for presentation use. It is different from the main report. The main report is for formal submission. This document is for \textbf{explaining the project clearly to the professor}, even if you do not have deep technical confidence.

Each major topic is explained in three layers:
\begin{itemize}
    \item simple meaning,
    \item technical meaning,
    \item what you can directly say in front of the professor.
\end{itemize}

\simple{You can use this file like a speaking guide.}

\technical{This guide translates the formal architecture, use cases, database design, and technology choices from the report into viva-ready language.}

% ---------------------------------------------------------------------------
\section{One-Minute Project Explanation}

\say{SecureGate is a smart access-control system for a campus or institution. Instead of using manual gate registers or easy-to-share ID cards, it gives users secure QR-based passes. Admins approve access, guards verify passes, contacts can see movement history, and the system maintains real-time logs. The completed version supports approval, QR verification, geofence-aware instant passes, optional face verification, and slot-based scheduling as the main course-level enhancement.}

If the professor wants a shorter version, say this:

\say{My project digitizes and secures the gate-entry process. It uses authentication, signed QR passes, optional face verification, GPS validation for instant passes, admin monitoring, and slot-based scheduling so entry and exit can be controlled by time window and capacity.}

% ---------------------------------------------------------------------------
\section{Problem Statement}

\subsection{What Problem Does This Project Solve?}

\simple{In many campuses, gate entry is still handled manually. That creates problems like slow verification, fake entries, missing logs, poor monitoring, and security gaps.}

\technical{Traditional gate control depends on handwritten registers, verbal approval, or easily shareable identity proof. These methods do not provide tamper resistance, real-time monitoring, policy enforcement, or auditable digital records.}

\say{The main problem is that manual gate management is slow, insecure, and hard to audit. SecureGate solves this by replacing manual access with digitally approved, verifiable, and logged access control.}

\subsection{Why Is This a Good Systems Project?}

\simple{Because it is not only a website. It is a complete system with different users, rules, data storage, real-time updates, security checks, and policy control.}

\technical{The project combines multi-role design, backend services, database persistence, cryptographic validation, real-time logging, optional face verification, geospatial checks, and deployment. This makes it a strong applied systems project rather than a simple CRUD application.}

\say{This is a systems project because it brings together authentication, secure pass generation, real-time event handling, database design, policy enforcement, and multi-portal coordination.}

% ---------------------------------------------------------------------------
\section{What SecureGate Actually Is}

\subsection{High-Level Idea}

\simple{SecureGate is a digital gate-pass platform. Different people use different portals. Members request access or passes. Admins approve and control policies. Guards verify passes. Contacts can view movement history.}

\technical{SecureGate is a role-based access-control platform implemented as a FastAPI-backed web system with separate frontend portals for members, admins, guards, and contacts. It manages the lifecycle of identity approval, pass creation, pass verification, and scan logging.}

\say{The simplest way to understand SecureGate is that it manages who is allowed to enter or exit, when they are allowed, and how that access is verified and recorded.}

\subsection{Who Uses the System?}

\begin{tabularx}{\textwidth}{@{}L{2.4cm}Y Y@{}}
\toprule
\textbf{Actor} & \textbf{Simple Meaning} & \textbf{Technical Role} \\
\midrule
Member & Regular system user who needs access & Authenticated user who requests access, books slots, or generates passes \\
Admin & Controls the whole system & Reviews requests, manages policies, configures geofence and slots, monitors logs \\
Guard & Verifies people at the gate & Uses scanner interface to validate QR, face match, expiry, and slot rules \\
Contact & Family or guardian type viewer & Accesses signed link to view recent movement and subscribe to alerts \\
Visitor & Guest user & Receives restricted one-time or scheduled pass issued by admin \\
\bottomrule
\end{tabularx}

\say{There are five main actors: member, admin, guard, contact, and visitor. Each actor has a different responsibility in the access-control process.}

% ---------------------------------------------------------------------------
\section{Full System Flow}

\subsection{End-to-End Story}

\begin{enumerate}
    \item A new member first requests account access.
    \item The admin reviews the request and approves or rejects it.
    \item After approval, the member can request a normal pass or use an instant pass flow.
    \item In the improved version, the member can also book a time slot.
    \item The system generates a secure QR token for approved access.
    \item At the gate, the guard scans the QR code.
    \item The backend checks whether the QR is genuine, valid, unused, and within policy.
    \item If needed, face verification is also performed.
    \item The result is stored in scan logs and reflected in the admin dashboard.
    \item Contacts can view recent movement through a secure signed link.
\end{enumerate}

\simple{The system is basically: request, approve, generate pass, scan, verify, log, monitor.}

\technical{This is a controlled credential lifecycle with approval workflow, secure token generation, runtime verification, and immutable audit recording.}

\say{The project flow is: user request, admin approval, QR generation, guard verification, and audit logging.}

% ---------------------------------------------------------------------------
\section{Architecture Explanation}

\subsection{What the Architecture Means}

\simple{The architecture is divided into layers so that each part has one job. Frontend portals collect input. The backend applies rules. The database stores records. External services handle notifications and maps.}

\technical{The project uses a layered design. The presentation layer contains the portals. The business logic layer contains FastAPI services such as authentication, pass handling, slot scheduling, face verification, geofence policy, and analytics. The persistence layer uses SQLAlchemy on SQLite or PostgreSQL. External services support notifications and map-based configuration.}

\say{I divided the system into presentation, backend logic, database, and integrations. This makes the system easier to manage and extend.}

\subsection{Layer-by-Layer Explanation}

\subsubsection{Presentation Layer}

\simple{This is the part users directly see.}

\technical{It contains separate interfaces for members, admins, guards, and contacts. This separation reduces confusion and lets each actor perform only the relevant actions.}

\say{The presentation layer is split by actor so each user sees only the features they need.}

\subsubsection{Business Logic Layer}

\simple{This is the brain of the project.}

\technical{The FastAPI backend handles authentication, approval workflow, pass creation, QR verification, slot management, geofence checks, face verification, notifications, and analytics.}

\say{The backend is where all the actual rules are applied. It decides whether a request is valid, whether a pass should be generated, and whether a guard should allow access.}

\subsubsection{Data Layer}

\simple{This is where the project stores everything permanently.}

\technical{The database stores users, registration requests, passes, scan logs, and in the improved design, slots and slot bookings. SQLAlchemy is used as the ORM to map Python models to database tables.}

\say{The data layer stores every important action so the system is auditable and not just temporary.}

\subsubsection{External Integrations}

\simple{These are supporting services that improve the system.}

\technical{Firebase FCM supports push notifications, Twilio supports SMS alerts, Leaflet is used for map-based geofence visualization, and Chart.js is used for analytics graphs.}

\say{External services are optional helpers. They are not the core logic, but they make the system more practical and complete.}

% ---------------------------------------------------------------------------
\section{Database and Data Model Explanation}

\subsection{Why the Database Matters}

\simple{The database is important because the system must remember who requested access, who got approved, which pass was used, and when someone entered or exited.}

\technical{The system depends on persistent state. Every approval, pass issuance, scan result, and policy-controlled action must be stored for consistency, auditing, and analytics.}

\say{Without the database, SecureGate would not be reliable, because the system must maintain history and verify state across multiple users and events.}

\subsection{Core Entities}

\footnotesize
\begin{tabularx}{\textwidth}{@{}L{3.25cm}Y@{}}
\toprule
\textbf{Entity} & \textbf{Explanation} \\
\midrule
\texttt{users} & Stores all active users, their role, credentials, face-related fields, and contact details \\
\texttt{registration\_requests} & Stores pending or reviewed account access requests before account creation \\
\texttt{passes} & Stores pass type, status, reason, QR token, validity, and usage state \\
\texttt{scan\_logs} & Stores every successful or failed checkpoint event for auditing and analytics \\
\texttt{slots} & Implemented entity for time windows, gate control, capacity, and policy flags \\
\shortstack[l]{\texttt{slot\_}\\\texttt{bookings}} & Implemented entity connecting users to reserved slots and tracking booking state \\
\bottomrule
\end{tabularx}
\normalsize

\say{The important idea is that each table represents one responsibility: identity, request review, access pass, audit log, and in the improved version, scheduling.}

% ---------------------------------------------------------------------------
\section{Security Explanation}

\subsection{Why Security Is Important Here}

\simple{Because this project decides who is allowed to enter or leave. If the system is weak, someone could fake a pass or misuse another person’s access.}

\technical{Access-control systems must prevent forgery, replay, unauthorized privilege use, and untracked movement. SecureGate addresses this through signed QR payloads, authentication, role checks, and audit logging.}

\say{Security is one of the main strengths of the project, because a gate-access system must not only work, it must also be trustworthy.}

\subsection{Main Security Features}

\subsubsection{JWT Authentication}

\simple{JWT means the system gives a secure login token after a user signs in.}

\technical{JWT is a signed token used to maintain authenticated sessions between frontend and backend without storing everything in server memory.}

\say{JWT helps the backend identify who is making a request after login.}

\subsubsection{Role-Based Access Control}

\simple{Different people should not be able to do the same things.}

\technical{RBAC enforces authorization by role. Admins can approve and configure. Guards can verify. Members can request. Contacts can only view what their signed link allows.}

\say{RBAC means permissions are controlled by role, not just by being logged in.}

\subsubsection{Signed QR Passes}

\simple{The QR code is not just plain text. It is protected so it cannot be easily changed or faked.}

\technical{The backend generates QR payloads with HMAC-SHA256 based signing. During scanning, the server recomputes the signature and rejects any tampered token.}

\say{The QR code is cryptographically signed, so if someone changes it, the system can detect that immediately.}

\subsubsection{Replay Prevention}

\simple{A pass that is already used should not work again.}

\technical{Pass state and scan logs are checked so one-time passes cannot be reused after successful verification.}

\say{Replay prevention means an old used pass cannot be reused to enter again.}

\subsubsection{Audit Logging}

\simple{Every important action is recorded.}

\technical{Every scan result and emergency event is stored in \texttt{scan\_logs}, enabling traceability, analytics, and compliance review.}

\say{Audit logging is important because security systems must leave evidence of what happened.}

% ---------------------------------------------------------------------------
\section{GPS and Geofence Explanation}

\subsection{What GPS Means in This Project}

\simple{GPS is used only in the instant pass feature. It checks whether the user is physically inside the allowed campus area before generating an immediate pass.}

\technical{The browser collects latitude and longitude. The backend calculates distance from the configured geofence center or boundary and compares it with the allowed radius or polygon settings.}

\say{GPS is not the whole project. It is one policy mechanism used only for instant on-campus pass generation and admin policy control.}

\subsection{Why the Professor Commented About GPS}

\simple{Because in the earlier write-up, GPS appeared too many times and made the use cases look repetitive.}

\technical{A strong use-case model should focus on actor goals, not repeat one technical mechanism across many sections. That is why the revised report limits GPS mainly to UC06 and UC11.}

\say{I improved the report by reducing GPS repetition and keeping it only where it is functionally important.}

% ---------------------------------------------------------------------------
\section{Slot-Based Scheduling Explanation}

\subsection{What Is a Slot?}

\simple{A slot means a fixed time window for entry or exit, like 9:00 to 9:15 AM, with a limited number of users.}

\technical{A slot is a policy-controlled scheduling unit defined by start time, end time, gate, capacity, and optional restrictions such as GPS or face verification requirements.}

\say{Slot scheduling is the main advanced enhancement that makes the system stronger at course level.}

\subsection{Why Add Slots?}

\simple{Because the professor wanted more distinct use cases and a stronger feature. Slots add time-based control and reduce crowding.}

\technical{Slots turn the system from a simple pass platform into a scheduling-aware access-governance platform. They introduce capacity control, temporal validation, and stronger policy enforcement.}

\say{Slots improve the project because they add scheduling, capacity planning, and rule-based verification, not just QR scanning.}

\subsection{How Slots Will Work}

\begin{enumerate}
    \item Admin creates slot definitions.
    \item Member chooses an available slot.
    \item System checks capacity and conflicts.
    \item Booking is stored.
    \item A pass request is linked to that slot.
    \item Guard verifies whether the pass is being used within the correct slot window.
\end{enumerate}

\say{The slot feature adds the answer to the professor’s question about scheduling and gives the system a stronger practical use in real campus management.}

% ---------------------------------------------------------------------------
\section{Explanation of the 12 Use Cases}

\subsection{How to Talk About Use Cases}

\simple{A use case means one meaningful goal that one actor performs in the system.}

\technical{Use cases should describe actor-system interaction around a complete goal, not just internal technical functions.}

\say{In the revised version, I made the use cases actor-centered so each one represents a different real activity in the system.}

\subsection{UC01 -- Access Request Submission}
\simple{A new user asks for access to use the system.}
\technical{The system validates details and stores a pending registration request for admin review.}
\say{This use case starts the user lifecycle.}

\subsection{UC02 -- Account Review and Provisioning}
\simple{The admin decides whether that new user should be allowed.}
\technical{The request is approved or rejected, and on approval, a real user account is created.}
\say{This ensures only verified users enter the system.}

\subsection{UC03 -- Slot Definition and Capacity Planning}
\simple{The admin defines available access windows.}
\technical{Slot objects include timing, gate association, capacity, and policy flags.}
\say{This is one of the new course-level additions.}

\subsection{UC04 -- Slot Booking by Member}
\simple{The member reserves a time window.}
\technical{The booking is recorded only if capacity and conflict checks pass.}
\say{This introduces scheduling control from the user side.}

\subsection{UC05 -- Standard Pass Request with Admin Approval}
\simple{The user requests a normal pass and waits for admin decision.}
\technical{A pass request is stored, reviewed, and converted to an HMAC-signed QR credential if approved.}
\say{This is the regular access process.}

\subsection{UC06 -- Instant On-Campus Pass Generation}
\simple{If the user is physically inside the campus, the system can generate a pass immediately.}
\technical{Browser geolocation is validated against the geofence before auto-approval.}
\say{This is where GPS is used in a focused and meaningful way.}

\subsection{UC07 -- Visitor Pass Creation}
\simple{An admin creates a pass for an external guest.}
\technical{A restricted access record is created with tighter validity rules and optional slot binding.}
\say{This supports guest handling without giving a full account.}

\subsection{UC08 -- QR Verification at the Checkpoint}
\simple{The guard scans the QR code and the system checks if it is valid.}
\technical{The backend validates signature, expiry, usage state, and optionally slot timing.}
\say{This is the core runtime verification process.}

\subsection{UC09 -- Restricted Access Face Verification}
\simple{For stronger checking, the guard can compare the person’s face with enrolled data.}
\technical{The system extracts live image features and compares them with stored templates.}
\say{Face verification is optional but adds a higher level of identity assurance.}

\subsection{UC10 -- Contact Subscription and History Review}
\simple{A contact can see the recent movement history of a member through a secure link.}
\technical{The signed token limits access to one mapped member without requiring a full contact account.}
\say{This adds accountability and transparency for guardians or family contacts.}

\subsection{UC11 -- Live Monitoring and Policy Adjustment}
\simple{The admin watches live gate activity and adjusts settings.}
\technical{Dashboard updates are driven by backend events, and policy values such as geofence or slots can be changed.}
\say{This makes the system operational, not static.}

\subsection{UC12 -- Emergency Exit Handling}
\simple{A user can leave urgently in an emergency.}
\technical{The event bypasses standard pass approval and is logged immediately for oversight.}
\say{Emergency handling shows the system can support real-world exceptions safely.}

% ---------------------------------------------------------------------------
\section{Technology Stack Explanation}

\subsection{How to Explain the Stack Simply}

\simple{The technology stack is just the set of tools used to build the system.}

\say{I selected each technology based on what job it performs in the system.}

\begin{tabularx}{\textwidth}{@{}L{3cm}Y@{}}
\toprule
\textbf{Technology} & \textbf{How to explain it} \\
\midrule
Python & Main programming language used to build the backend logic \\
FastAPI & Framework used to create the backend API quickly and clearly \\
Uvicorn & Server that runs the FastAPI application \\
SQLAlchemy & Tool that connects Python code to the database \\
Pydantic & Tool used to validate request and response data \\
PyJWT & Used for login tokens and secure signed access links \\
Passlib + bcrypt & Used for password hashing so passwords are not stored as plain text \\
HMAC-SHA256 & Used to sign QR tokens so they cannot be forged easily \\
OpenCV & Used for image handling and face-related operations \\
face-recognition & Optional heavier face matching library when enabled \\
geopy / Shapely & Used for geofence and location logic \\
Firebase / Twilio & Used for notifications like push alerts and SMS \\
jsQR / QRCode.js & Used on the frontend to scan and generate QR codes \\
Chart.js / Leaflet & Used for graphs and map-based admin controls \\
\bottomrule
\end{tabularx}

\subsection{How to Answer ``Why These Technologies?''}

\say{I chose FastAPI because it is clean and efficient for backend APIs, SQLAlchemy because it helps manage database models properly, JWT for secure session handling, HMAC for QR integrity, OpenCV for vision-related tasks, and geopy or Shapely for location enforcement.}

% ---------------------------------------------------------------------------
\section{What Is Already Implemented and What Is Planned}

\subsection{Already Implemented}

\begin{itemize}
    \item account request and admin approval,
    \item authentication and role control,
    \item signed QR pass generation,
    \item guard scanning and verification,
    \item geofence-aware instant pass generation,
    \item optional face registration and matching,
    \item real-time scan logs and dashboard,
    \item contact history view through signed link,
    \item emergency exit handling,
    \item deployed version on Render.
\end{itemize}

\subsection{Planned as the Main Course-Level Extension}

\begin{itemize}
    \item slot creation,
    \item slot booking,
    \item capacity-aware scheduling,
    \item slot-window enforcement during verification,
    \item stronger analytics and test coverage.
\end{itemize}

\say{I am careful to separate what is already deployed from what is the next planned extension, so the project remains honest and professional.}

% ---------------------------------------------------------------------------
\section{Likely Professor Questions and Good Answers}

\subsection{Why is this not just a normal website?}
\say{Because it is not only displaying information. It manages secure credentials, approval workflows, verification at runtime, policy enforcement, logging, and real-time monitoring across multiple user roles.}

\subsection{Why did you add slot scheduling?}
\say{Slot scheduling adds time-window control and capacity management, which makes the system stronger, more practical, and more suitable as a course-level systems enhancement.}

\subsection{Why use QR instead of only face recognition?}
\say{QR is faster and cheaper for routine verification, while face recognition can be used only when stronger identity assurance is needed. That gives a better balance between speed and security.}

\subsection{Why is GPS not a separate major module everywhere?}
\say{Because GPS is only one policy mechanism. In the revised report, I limited it to the instant pass flow and admin policy configuration so the use cases remain distinct and not repetitive.}

\subsection{How is the QR code secured?}
\say{The QR token is signed using HMAC-SHA256, so if someone modifies the token, the system detects the mismatch during verification.}

\subsection{What makes this project academically strong?}
\say{It combines secure access control, role-based design, auditability, policy-driven validation, optional biometrics, geospatial checks, real-time monitoring, and a clear extension path through slot scheduling.}

\subsection{What is your main contribution in the revised report?}
\say{My main contribution is restructuring the project into 12 clear actor-centered use cases and introducing slot-based scheduling as the major course-level enhancement.}

% ---------------------------------------------------------------------------
\section{How to Present the Architecture Diagram}

\say{At the top, I have the user-facing portals. In the middle, I have the backend services such as authentication, pass handling, slot scheduling, face verification, geofence policy, notifications, and analytics. Below that, I have the database layer where core entities are stored. Finally, I have external services like Firebase, Twilio, Leaflet, and chart libraries. This layered structure makes the system modular and easier to maintain.}

\simple{When explaining the architecture, do not rush into package names first. Start with the job of each layer.}

% ---------------------------------------------------------------------------
\section{How to Present the Use Cases}

\say{Earlier the use cases were too similar because some technical features like GPS appeared repeatedly. In the revised version, every use case now represents one distinct actor goal: registration, approval, slot definition, slot booking, normal pass request, instant pass, visitor pass, QR verification, face verification, contact history, live monitoring, and emergency exit.}

\simple{This is a strong answer if the professor asks why the use cases changed.}

% ---------------------------------------------------------------------------
\section{How to Conclude the Project}

\say{SecureGate started as a working smart gate-pass system, and the revised course-level version makes it more structured and stronger by introducing clear actor-based use cases, policy-focused design, and implemented slot-based scheduling. The system now proves the main workflow plus time-based control, capacity management, analytics, and stronger testing.}

% ---------------------------------------------------------------------------
\section{Final Speaking Tips}

\begin{itemize}
    \item Start from the problem, not from the technology.
    \item Say what is implemented and what remains optional or future work.
    \item Use the words \textbf{secure}, \textbf{policy-driven}, \textbf{auditable}, and \textbf{role-based}.
    \item If you forget a term, explain the job instead of the package name.
    \item If the professor asks something deep, first answer simply, then add the technical meaning.
    \item Do not say GPS is the whole project. Say it is one part of policy validation.
    \item Say slots are implemented for booking, pass binding, and guard time-window enforcement.
\end{itemize}

\section{Final Closing Line}

\say{SecureGate is a secure, role-based, and auditable access-control platform that manages the full lifecycle of campus entry and exit. Its course-level strength comes from combining digital approval, cryptographically signed QR passes, optional face verification, geofence-aware validation, real-time monitoring, and implemented slot-based scheduling.}

\end{document}
